% Block A Methods. Working draft: analysis/block_a/METHODS_DRAFT.md
% [PI] marks a decision recorded in analysis/block_a/DISCREPANCY_REPORT.md.

\section*{Methods}

\subsection*{Panel}

Predictions were run for 48 G-protein-coupled receptors --- 40 Class~A, 4
Class~B and 4 Class~F --- under two input conditions and on four co-folding
backbones, giving 380 receptor $\times$ backbone $\times$ arm cells and 9{,}490
scored predictions. Cells hold a median of 25 predictions (minimum 20; the
Boltz-2 cognate arm is short 10 rows, a documented gap).

Each receptor is paired with deposited active- and inactive-state reference
structures. The reference set comprises \textbf{[PI: 98 empirical / 89 as
previously stated]} unique Protein Data Bank entries across the panel --- 44 in
an active-role and 54 in an inactive-role assignment --- drawn from a total
reference collection of 167 entries. The three denominators are reconciled in
Table~S3. This denominator must be stated wherever the reference set is
enumerated, because the panel (48), the panel-PDB set (98) and the full
collection (167) are all defensible counts of different things.

\subsection*{Conformational predicates}

Two intrinsic geometric measurements are computed on each predicted structure
alone. Neither is a comparison to a reference.

\texttt{d\_npxxy\_oh} is the hydroxyl--hydroxyl distance between Tyr5.58 and
Tyr7.53. On activation Tyr7.53 of the NPxxY motif rotates into the receptor core
and packs against Tyr5.58 in an arrangement that is usually water-mediated. A
prediction is called active on this axis below 9.080~\AA.

\texttt{d\_gpcrdb\_tm6\_tilt\_246\_637\_ca} is a C$\alpha$--C$\alpha$ distance
capturing outward displacement of the cytoplasmic end of TM6, the canonical
activation signature \citep{hilger2020gcgr}. A prediction is called active above 14.932~\AA.

A prediction is called active only when both fire. Thresholds are carried as
row-level columns and were read from the data rather than reimplemented. Both
were derived from 80 tier-1 crystallographic reference rows by a midpoint
objective, and from no prediction. \textbf{Self-classification accuracy on that
set is a derivation-set figure, not a held-out one}, and is reported as such.

The thresholds are derived from geometry rather than fitted to a classifier, but
the 80 rows carry state annotations, and the provenance of those annotations
bounds how independent the instrument is: an instrument calibrated on a given
labelling cannot then adjudicate that labelling.

This is not a hypothetical concern. Curated activation-state labels are not
purely geometric --- they are assigned from ``GPCRdb annotations, structural
criteria, and ligand binding status'' \citep[p.~12]{khaleq2026hyaline} --- and
an unsupervised geometric index over 1{,}351 class~A structures recovers
entries ``whose state of activity was originally interpreted incorrectly''
\citep[p.~3]{paajanen2026activation}. The same work notes that many models
predate substantial active-state coverage, so ``it is also not certain that the
active states of the structures used in developing the models were correctly
identified'' (p.~3).

Our predicate follows the geometric pattern rather than the label-inheriting
one: it reads C$\alpha$ and hydroxyl positions at Ballesteros--Weinstein
positions 3.50, 5.58, 6.34 and 7.53, using generic \emph{numbering} and not
curated \emph{labels}. The residual exposure is the selection of the 80 tier-1
rows themselves. \textbf{[PI: were those rows selected by crystallographic tier,
with labels checked only afterwards, or selected by curated label?]} The first
leaves the instrument independent of the annotation; the second does not, and
would need saying. The same reference set's construct annotation contradicts the
RCSB record on roughly 40\% of evaluable entries, so its labelling is not
assumed reliable elsewhere in this work.
The models do not place water molecules; results are described as consistent
with the water-bridged arrangement, never as the bridge forming.

\subsection*{Prediction protocol}

Four co-folding models were used: Boltz-2, Chai-1, OpenFold3 and Protenix2. Each
received either the receptor sequence alone (apo arm) or the receptor sequence
together with its cognate G$\alpha$ subunit (cognate arm). Templates were
disabled on all backbones, multiple sequence alignments were pinned, and no
reference structure was supplied to any model at any point. Twenty-five seeds
were run per condition.

The apo arm is a computational reference condition, not a physical state --- no
ligand, no membrane, and no basal-activity equivalent. Real receptors show
constitutive activity that varies widely between receptors, and their allosteric
coupling differs correspondingly \citep{georgiou2025heterogeneity}; the apo arm
should not be equated with a physiological inactive state.

OpenFold3 seeds were generated as $5 \times 5$ outer/inner after a seeding fix,
so OF3 rows are less independent than the other backbones'. This qualifies any
cross-backbone comparison of correlation magnitudes. Scoring used a single
frozen scorer (version 0.1.0, commit \texttt{04243c45}) recorded per row.

\subsection*{Uncertainty}

The cluster bootstrap is authoritative: 1{,}000 resamples over 26 paralog
clusters, seed 20260909. A receptor-level bootstrap is reported as secondary, in
supplementary tables only. The cluster interval is equal to or wider than the
receptor interval on every headline statistic, by a factor of up to 2.2.

Eleven of the 26 clusters are singletons (42\%), so the paralogy correction acts
on only 15 multi-member clusters. All confidence intervals quoted in the main
text and in Tables~2--4 are cluster intervals.

\subsection*{Exclusions}

Five exclusion sets were pre-specified by cause and applied before results were
inspected. They are not a quality filter and are not applied as a union.

E1 (25 rows) removes one cell whose mean pLDDT falls below 50. E2 (4 rows)
removes model-side atom clashes with \texttt{d\_npxxy\_oh} below 2.4~\AA. These
two together remove 29 rows (0.31\%) and are applied everywhere.

E3 (4{,}890 rows) flags receptors whose reference fails its own predicate on a
given axis. \textbf{It is applied per axis} --- \texttt{excl\_E3\_npxxy}
(4{,}690 rows, 24 receptors) or \texttt{excl\_E3\_tilt} (2{,}000 rows, 10
receptors) --- and only where a reference value serves as a denominator or a
regression predictor. It merges two distinct causes with different
consequences, and we report them separately. On the NPxxY axis, 12 receptors
(2{,}295 rows) have no measurable value because the axis is undefined for them
--- EDNRB, for example, carries Leu at 7.53, so an NPxxY hydroxyl distance does
not exist --- while a further 12 receptors (2{,}395 rows) have a measurable
reference that fails the predicate. The first is missing data; the second is an
observation about reference quality. On the tilt axis all 10 flagged receptors
are of the second kind and none is unmeasurable. The E3 union is never applied, and E3 is irrelevant to raw
distributions, predicate firing rates and confidence correlations, which never
touch a reference separation.

E4 (1{,}495 rows) restricts to Class~A; Class~B and~F are reported separately as
instrument scope. E5 (500 rows) marks agonist-only actives and is shown as a
sensitivity contrast, never applied silently.

Every headline statistic was recomputed under every exclusion combination
(Table~5). No sign flip occurs anywhere and no statistic moves by more than
0.5\% from baseline.

\subsection*{Confidence aggregation}

Three aggregations of pLDDT were computed: the whole-complex mean, the mean over
the six state-defining anchor residues 3.50, 3.51, 5.58, 6.30, 6.34 and 7.53,
and the minimum over those anchors. Correlations are computed over approximately
1{,}600 rows and 32 receptors per backbone.

\textbf{The anchor-restricted mean was designated the primary aggregation post
hoc, after all three had been computed.} All three are reported in Table~4. The
disclosure, not the designation, is what the result rests on.

\subsection*{Data-layer verification}

Both predicate axes were independently recomputed against non-scorer
implementations and agree bit-exactly (Pearson $r = 1.000000$). The predicate
booleans were recomputed from the raw axis values rather than trusted, and cell
counts were reconciled against the run manifest. Fold integrity was checked
independently: the TM6 helicity anchor holds on 96.3\% of rows.

\subsection*{Known data-layer issues}

Four rows carry sub-covalent hydroxyl--hydroxyl distances, which are model-side
clashes rather than measurements, and are excluded as E2. One 25-row cell passed
every assertion in the verification suite because the suite has no pLDDT floor,
and is excluded as E1. The NPxxY threshold was truncated to 9.08 from 9.082,
which changes the call for one borderline row.

\subsection*{Reference-set limitations}

The construct annotation carried with the reference set contradicts the RCSB
record on approximately 40\% of evaluable entries; the denominator for that rate
is \textbf{[PI: 162 empirical / 127 as previously stated]}. Several references
carry fusion partners or engineered residues inside the predicate windows, which
is recorded per entry in Table~S1. Nine references deviate from their expected
predicate call; five of the nine are unclassified as to cause, with the
remainder split between expected biology, curation error, a measurement
artifact, and one entry that could be either. Several entries annotated active
were solved with agonist bound but no transducer and do not form the NPxxY
network, which the predicate reports correctly.
